% discussion.tex — Section IV: Discussion
% Target: PRD two-column (revtex4-2)
% Macros used: \HH, \dL, \Mz, \Mbh, \Mstellar, \pGW, \pdet, \Lcat, \Lcomp, \SNR, \pending, \todo

The results of Sec.~\ref{sec:results} demonstrate that the redshifted MBH
mass $\Mz$ provides a powerful additional handle for dark siren cosmology
with EMRIs.  We now place these findings in the context of prior work,
discuss the physical mechanism behind the improvement, and assess the
limitations of the analysis.

%% ============================================================
\subsection{Comparison with previous work}
\label{sec:comparison}
%% ============================================================

The most direct comparison is with
Laghi~\emph{et~al.}~\cite{Laghi:2021pqk}, who performed the first
comprehensive study of EMRI dark sirens as probes of $\HH$ with LISA.
Our analysis extends their work in four significant ways.

First, Laghi~\emph{et~al.}\ adopted a fixed fractional luminosity
distance uncertainty of $\sigma_{\dL}/\dL = 10\%$ for all detected
EMRIs, independent of source parameters.  We replace this simplification
with the full Cram\'{e}r--Rao bound derived from the $14 \times 14$
Fisher information matrix (Sec.~\ref{sec:fisher}), which captures the
dependence of $\sigma_{\dL}$ on SNR, sky position, mass ratio, and spin.
The resulting per-event distance uncertainties span a wide range, with
some high-SNR detections achieving $\sigma_{\dL}/\dL < 1\%$.

Second, we introduce the MBH mass channel, which exploits the precise
measurement of $\Mz$ from the gravitational waveform as an additional
discriminant for host galaxy identification.  This channel is absent from
Ref.~\cite{Laghi:2021pqk} and is, to our knowledge, the first
application of the MBH mass information to EMRI dark siren cosmology.

Third, we employ the real GLADE+ galaxy catalog~\cite{Dalya:2021ewn},
which includes spectroscopic and photometric redshifts, stellar masses,
and sky positions for approximately 23~million objects.
Laghi~\emph{et~al.}\ used a synthetic catalog with idealized properties.
The real catalog introduces the practical complication of spatially
varying completeness, which we address by implementing the correction of
Gray~\emph{et~al.}~\cite{Gray:2019ksv}.

Fourth, we include the catalog completeness correction
[Eq.~\eqref{eq:completeness_combination}] following Gray~\emph{et~al.},
which was originally developed for ground-based dark siren analyses of
binary black hole mergers~\cite{Gray:2019ksv}.  This correction is
essential for avoiding the systematic low-$\HH$ bias caused by
preferentially cataloged low-redshift galaxies, as demonstrated in
Sec.~\ref{sec:systematics}.

Our Fisher matrix results are validated against the benchmark
Cram\'{e}r--Rao bounds in Table~III of
Babak~\emph{et~al.}~\cite{Babak:2017tow}, finding agreement to within
the numerical precision of the finite-difference stencil
(Sec.~\ref{sec:pe_validation}).  This provides confidence that the
parameter estimation inputs to the dark siren likelihood are correctly
computed.

The completeness correction we implement follows the framework of
Gray~\emph{et~al.}~\cite{Gray:2019ksv}, adapted here for EMRI sources
with LISA.  Gair~\emph{et~al.}~\cite{Gair:2022zsa} recently provided a
comprehensive study of systematic biases in dark siren analyses,
identifying catalog incompleteness as the dominant source of error for
events beyond $z \sim 0.03$ in GLADE+.  Our treatment of the
completeness correction is consistent with their recommendations,
although we note that our completeness function is interpolated from
digitized data rather than derived from a self-consistent luminosity
function model (see Sec.~\ref{sec:caveats}).

%% ============================================================
\subsection{Physical interpretation of the MBH mass improvement}
\label{sec:mass_interpretation}
%% ============================================================

The factor of $\sim 2$ improvement in $\HH$ precision from the MBH mass
channel (from $\sim 4\%$ to $\sim 1.8\%$ in the $\pdet = 1$ baseline)
has a clear physical origin: the mass measurement breaks the
luminosity-distance--redshift degeneracy inherent in dark siren analyses.

Without the mass channel, a GW detection constrains $\dL$ (and sky
position) but not $z$ independently.  For a given trial value of $\HH$,
the inferred redshift $z = z(\dL, \HH)$ identifies a shell in the galaxy
catalog.  All galaxies consistent with the measured distance and sky
position contribute to the single-event likelihood, potentially
numbering $\sim 10$--$10^4$ candidates depending on the sky
localization volume and local galaxy density.  Because each candidate is
weighted equally (up to the GW measurement uncertainty), the per-event
$\HH$ constraint is diluted by the large number of viable hosts.

With the mass channel, the precisely measured $\Mz$ imposes an
additional constraint: the true host must satisfy
$\Mbh \approx \Mz / (1+z)$, where $z$ is determined by $\dL$ and the
trial $\HH$.  Since the GW measurement determines $\Mz$ to a fractional
precision of $\lesssim 10^{-6}$ (Sec.~\ref{sec:pe_validation}), while
the catalog masses inferred from the $\Mbh$--$\Mstellar$ scaling
relation~\cite{Reines:2015moa} carry uncertainties of $\sim 20\%$, the
mass match effectively acts as a narrow filter.  Only galaxies whose
inferred BH mass is consistent with $\Mz / (1+z)$ within the scaling
relation scatter contribute significantly to the likelihood
[Eq.~\eqref{eq:Lcat_mass}].  This typically reduces the number of viable
hosts by one to two orders of magnitude, concentrating the single-event
likelihood around the correct redshift and thus the correct $\HH$.

The improvement is largest for detections in regions of high galaxy
density (large angular extent and moderate distance), where the standard
channel admits the most candidate hosts.  Conversely, for isolated
detections in underdense regions where only a handful of galaxies lie
within the error volume, the mass information provides a smaller
marginal improvement because the host galaxy is already effectively
identified by proximity.

We note that the efficacy of the mass channel is limited by the accuracy
of the $\Mbh$--$\Mstellar$ scaling relation.  The
Reines~\&~Volonteri~\cite{Reines:2015moa} relation has an intrinsic
scatter of approximately $0.55\,\mathrm{dex}$, with best-fit parameters
$\alpha = 7.45 \pm 0.08$ and $\beta = 1.05 \pm 0.11$.  Improvements in
the empirical calibration of this relation, particularly in the
$10^5$--$10^6\,M_\odot$ range relevant for EMRI sources, would directly
translate into a tighter mass filter and thus a stronger $\HH$
constraint.

%% ============================================================
\subsection{Limitations and caveats}
\label{sec:caveats}
%% ============================================================

We identify several approximations and limitations of the present
analysis.  We discuss them in roughly decreasing order of expected
impact on the results.

\paragraph{Fisher matrix approximation.}
The parameter estimation relies on the Fisher information matrix, which
provides a lower bound on parameter variances that is tight only in the
high-SNR, unimodal posterior regime~\cite{Vallisneri:2007ev}.  While our
detection threshold of $\SNR \geq 20$ places the analysis in a regime
where the Fisher approximation is generally reliable, multimodal
posteriors can arise for EMRI sources due to correlations between mass,
spin, and orbital parameters.  In such cases, the Fisher matrix may
underestimate the true parameter uncertainty, leading to overly
optimistic Cram\'{e}r--Rao bounds.  A full Bayesian parameter estimation
using, e.g., Markov chain Monte Carlo (MCMC) sampling would capture
these effects but is computationally prohibitive for the $\sim 10^5$
injections considered here.

\paragraph{Finite-difference accuracy.}
The waveform derivatives entering the Fisher matrix are computed using
forward finite differences with step size $\epsilon = 10^{-6}$,
yielding $\mathcal{O}(\epsilon)$ accuracy.  The five-point stencil of
Eq.~\eqref{eq:five_point}, which achieves $\mathcal{O}(\epsilon^4)$
accuracy, is implemented in the codebase but was not used for the
injection campaign due to the $4\times$ increase in waveform evaluations
per parameter (56 vs.\ 14 per Fisher matrix).  While the forward-difference
and five-point results agree for the validation source of
Ref.~\cite{Babak:2017tow}, the higher-order stencil would provide
improved robustness for sources near the edges of parameter space.

\paragraph{Simplified population model.}
We draw EMRI source parameters from a random distribution over the
parameter ranges listed in Appendix~\ref{app:parameters}, rather than
from an astrophysically motivated population model.  Realistic EMRI rate
estimates~\cite{Babak:2017tow} predict event rates that depend strongly
on the MBH mass function, the stellar environment of the galactic
nucleus, and the efficiency of orbital capture.  A population-informed
injection campaign may yield a different detection yield and a different
distribution of per-event constraints, which could affect both the
overall $\HH$ precision and the relative improvement from the mass
channel.

\paragraph{$\Mbh$--$\Mstellar$ scaling relation.}
The mass channel relies on the empirical scaling relation of
Reines~\&~Volonteri~\cite{Reines:2015moa} to infer the MBH mass from
the stellar mass of candidate host galaxies.  This relation was
calibrated in the local universe ($z \lesssim 0.055$) and may evolve
with redshift.  Moreover, the intrinsic scatter of the relation
($\sim 0.55\,\mathrm{dex}$) is the dominant source of uncertainty in
the galaxy-side mass estimate.  Systematic shifts in the scaling
relation parameters $\alpha$ and $\beta$ would propagate directly into
the mass filter, potentially biasing or broadening the $\HH$ posterior.
Future work should marginalize over the scaling relation parameters
within their uncertainties.

\paragraph{GLADE+ catalog completeness.}
The completeness fraction $f(z, h)$ entering the correction of
Eq.~\eqref{eq:completeness_combination} is interpolated from published
completeness data for the GLADE+ catalog~\cite{Dalya:2021ewn}.  We note
two caveats.  First, the published completeness may reflect
\emph{number} completeness (fraction of galaxies detected) rather than
\emph{luminosity} completeness (fraction of total luminosity captured),
which is the more physically relevant quantity for galaxy-catalog-based
analyses~\cite{Gray:2019ksv}.  Second, the completeness is treated as a
function of redshift alone, whereas the actual catalog completeness
varies with sky position, galaxy type, and apparent magnitude.  A more
refined treatment would use the angular-dependent completeness, though
this is beyond the scope of the present work.

\paragraph{Cosmological model.}
We assume a flat $\Lambda$CDM cosmology with fixed matter density
$\Omega_m = 0.25$.  This value reflects the WMAP-era best
fit~\cite{Hinshaw:2012aka} rather than the Planck~2018 result
$\Omega_m = 0.3153 \pm 0.0073$~\cite{Planck:2018vyg}.  While the
choice of $\Omega_m$ affects the $\dL$--$z$ relation and hence the
inferred $\HH$, the sensitivity of the posterior to $\Omega_m$ is
subdominant to the statistical uncertainty from the finite number of
detections.  A joint inference over $(\HH, \Omega_m)$ or an extension
to the $w$CDM model is straightforward with the existing pipeline
infrastructure but is deferred to future work.

\paragraph{Galactic confusion noise.}
The LISA noise PSD used in this analysis
[Eq.~\eqref{eq:psd}] does not include the foreground from unresolved
galactic white-dwarf binaries, which dominates the sensitivity in the
$0.1$--$3\,\mathrm{mHz}$ band.  While the noise model includes the
$S_c(f)$ term in the formal expression, this component is set to zero in
the current implementation.  Including the confusion noise would reduce
the effective SNR for EMRI sources with significant power in this band,
potentially lowering the detection yield and increasing the per-event
distance uncertainty.  The effect is expected to be moderate for the MBH
mass range considered here ($10^5$--$10^6\,M_\odot$), whose peak GW
emission lies at frequencies above the confusion noise band.

%% ============================================================
\subsection{Future directions}
\label{sec:future}
%% ============================================================

Several extensions of this work are planned or in progress.

\paragraph{Completeness-corrected production run.}
The results presented in Sec.~\ref{sec:results} use the $\pdet = 1$
baseline.  A production-scale simulation incorporating the
completeness-corrected likelihood of
Eq.~\eqref{eq:completeness_combination} with a realistic $\pdet$
estimated from the injection campaign is underway.  This will quantify
the bias reduction from the Gray~\emph{et~al.}\ correction and
establish the definitive $\HH$ constraint.

\paragraph{Astrophysical population model.}
Replacing the uniform parameter draws with an astrophysically motivated
EMRI population model~\cite{Babak:2017tow} --- incorporating the MBH
mass function, stellar cusp properties, and orbital capture
efficiencies --- would provide a more realistic estimate of both the
detection yield over the LISA mission lifetime and the expected $\HH$
precision.

\paragraph{Full Bayesian parameter estimation.}
MCMC-based parameter estimation for individual EMRI detections would
replace the Fisher matrix approximation and correctly capture multimodal
and non-Gaussian posterior features.  While this is computationally
expensive for the full injection campaign, targeted MCMC runs on a
representative subset of detections would validate the Fisher-matrix
results and quantify any systematic bias from the Gaussian
approximation.

\paragraph{Five-point stencil derivatives.}
Upgrading the waveform derivatives from forward differences to the
five-point stencil of Eq.~\eqref{eq:five_point} for the full injection
campaign would improve the Fisher matrix accuracy at the cost of a
$4\times$ increase in waveform evaluations.  This is straightforward to
implement (the five-point stencil is already available in the codebase)
and would be worthwhile for a production-quality analysis.

\paragraph{Joint analysis with other LISA sources.}
EMRIs are one of several source classes observable by LISA.  Combining
the EMRI dark siren constraints with those from massive black hole
binaries (MBBHs) and stellar-origin binary black holes (SOBBHs) would
improve the overall $\HH$ precision, particularly if the different
source classes probe complementary redshift
ranges~\cite{LISA:2024hlh,Chen:2017rfc}.

\paragraph{Dark energy equation of state.}
The pipeline infrastructure already supports a $w$CDM cosmology with
parameters $(w_0, w_a)$, though this extension has not yet been
exercised.  With a sufficiently large detection catalog, EMRI dark
sirens could provide independent constraints on the dark energy equation
of state, complementary to those from Type~Ia
supernovae~\cite{DES:2024tys} and baryon acoustic
oscillations.

\paragraph{Improved completeness modeling.}
An angular-dependent and luminosity-weighted completeness model for
GLADE+ would improve the fidelity of the correction in
Eq.~\eqref{eq:completeness_combination}.  In combination with upcoming
spectroscopic surveys (e.g., DESI, Euclid), which will substantially
increase the catalog depth at $z \gtrsim 0.1$, the completeness
correction is expected to become less critical as the catalog approaches
full coverage over the EMRI detection volume.

%% CITATIONS NEEDED (for gpd-bibliographer):
%% - MISSING:wmap-9yr-2013 — Hinshaw et al. (2013), WMAP 9-year cosmological parameters, arXiv:1212.5226
%% - MISSING:des-y5-2024 — DES Collaboration, Year 5 supernova cosmology results
